\documentclass[11pt,a4paper]{article}
\usepackage[hmargin=20mm,vmargin=17mm,headheight=14pt,headsep=7mm,footskip=11mm]{geometry}
\usepackage[T1]{fontenc}
\usepackage{amsmath,amsthm}
\usepackage{newtxtext,newtxmath}
\usepackage{microtype}
\usepackage{booktabs,tabularx,array}
\usepackage{enumitem}
\usepackage{titlesec}
\usepackage{needspace}
\usepackage{xcolor}
\usepackage{listings}
\usepackage{fancyhdr}
\usepackage{hyperref}
\definecolor{ink}{HTML}{173A54}
\definecolor{codebg}{HTML}{F3F5F7}
\hypersetup{colorlinks=true,linkcolor=ink,citecolor=ink,urlcolor=ink,
  pdftitle={Exact Evaluation of QOp with Cyclotomic Arithmetic},
  pdfsubject={Exact evaluation prototype and student projects for QPP}}
\pagestyle{fancy}
\fancyhf{}
\fancyhead[L]{\small Exact evaluation of QOp}
\fancyhead[R]{\small QPP / HQP}
\fancyfoot[C]{\thepage}
\renewcommand{\headrulewidth}{0.3pt}
\setlength{\parindent}{0pt}
\setlength{\parskip}{2pt plus 1pt}
\titlespacing*{\section}{0pt}{2ex plus .5ex minus .2ex}{1.2ex plus .2ex}
\titlespacing*{\subsection}{0pt}{1.5ex plus .5ex minus .2ex}{1ex plus .2ex}
\titlespacing*{\paragraph}{0pt}{1.25ex plus .5ex minus .2ex}{1em}
\setlist{nosep,leftmargin=1.5em}
\setlength{\emergencystretch}{2em}
\lstset{language=Haskell,basicstyle=\fontsize{10}{11}\selectfont\ttfamily,
  backgroundcolor=\color{codebg},frame=single,rulecolor=\color{black!15},
  framerule=0.3pt,framesep=5pt,columns=fullflexible,
  keepspaces=true,showstringspaces=false,breaklines=true,
  aboveskip=5pt,belowskip=5pt,
  literate={⊗}{{$\otimes$}}1 {∘}{{$\circ$}}1}
\newcommand{\Q}{\mathbb{Q}}
\newcommand{\Z}{\mathbb{Z}}
\newcommand{\C}{\mathbb{C}}
\newcommand{\K}{K_N}
\newcommand{\D}{D_N}
\newcommand{\sem}[1]{\llbracket #1\rrbracket}
\newcommand{\code}[1]{\texttt{#1}}
\newcommand{\ket}[1]{|#1\rangle}
\DeclareMathOperator{\lcm}{lcm}
\DeclareMathOperator{\rank}{rank}
\DeclareMathOperator{\diag}{diag}
\DeclareMathOperator{\tr}{tr}
\newtheorem{proposition}{Proposition}
\newcolumntype{Y}{>{\raggedright\arraybackslash}X}

\begin{document}
\begin{center}
{\LARGE\bfseries Exact Evaluation of QOp\\[3pt] with Cyclotomic Arithmetic}\\[8pt]
{\large Mathematical foundations, two evaluators, and exact MPS projects}\\[7pt]
{\small QPP / HQP \quad$\cdot$\quad 17 October 2026}
\end{center}
\vspace{-5pt}
\begin{abstract}
\noindent Quantum Programming Playground (QPP) represents quantum operators
as symbolic programs whose angles are rational multiples of $\pi$.
Every valid finite expression has its entries in a computable cyclotomic
field. We derive this closure result and implement dense matrix and direct
statevector evaluation as two realizations of one semantics. Worked examples
connect syntax, coefficients, amplitudes, and probabilities. Exact matrix
product states need a different factorization: their amplitudes can remain
cyclotomic when their singular values cannot. We develop rank factorization
into a four-week HQP project and an MSc thesis proposal, with an independent
project on scalar efficiency. The prototype implements scalar arithmetic
and both evaluators, including measurement and reset through unnormalized
branches, with expected-value and randomized tests.
MPS remains proposed; no comparative performance claim is made.
\end{abstract}

\section{Why exact evaluation, and what must change?}
\label{sec:intro}
Compiler validation asks whether a rewrite preserves every amplitude or
an apparently zero entry actually vanishes. Exact evaluation decides these
questions without a tolerance. QPP is a Haskell framework whose symbolic
operator language, \code{QOp}, has several evaluators. Its \code{Rational}
angles become complex coefficients during evaluation: rationals alone
cannot represent $1/\sqrt2$ or a nonreal phase. Cyclotomic fields contain
both and support gate composition. Their degree and coefficient sizes
determine the arithmetic cost.

\paragraph{Scope and reading guide.}
We assume linear algebra and qubits, but no QPP or algebraic-number background.
Language, algebra, and examples lead into shared foundations, two evaluators,
costs, and MPS projects. Listings assume basic Haskell; equations specify their
meaning independently. Optional Galois arguments prove the limitations of SVD
and normalization. The prototype supports cyclotomic vectors, measurement,
and reset, but excludes arbitrary normalization.

\section{The operator language and its mathematical objects}
\label{sec:language}
A \code{QOp} constructor builds syntax, performing no numerical evaluation.

\paragraph{Terms, operators, and vectors.}
Constructor names appear in monospace; $u,t$ are terms. For a valid term
$u$, its width $n_u$ is the number of qubits on which it acts. Its
denotation $\sem u$ is a linear operator on $\C^{2^{n_u}}$, represented
in the computational basis by a $2^{n_u}\times2^{n_u}$ matrix.
The notation $\sem u\,v$ is the action of this same operator on a column
vector $v$. We keep these brackets for both evaluators.
Symbols such as $A=\sem u$ name mathematical matrices, not syntax.

The ordered basis is $\ket{b_0\cdots b_{n-1}}$, with
$b_j\in\{0,1\}$ and index $\sum_j b_j2^{n-1-j}$. Thus wire zero is
the most significant bit. For two qubits the vector order is
$(v_{00},v_{01},v_{10},v_{11})^T$. The left tensor factor acts on the
leading wires. A zero-qubit space has dimension one: \code{Id 0} denotes
$[1]$, and \code{Phase q} denotes $[e^{i\pi q}]$.

\paragraph{Widths and validity.}
\code{Compose u t} applies $t$ first and requires $n_u=n_t$.
\code{Tensor u t} has width $n_u+n_t$.
\code{DirectSum u t} requires $n_u=n_t$ and adds a selector qubit:
its width is $n_u+1$. Selector value zero chooses $u$ and value one
chooses $t$. In particular, \code{C u} has width $n_u+1$ and selects
between identity and $u$. Here \code{Nat} is represented by \code{Int},
so nonnegative widths must be checked explicitly.
A wire permutation must be a bijection of its local indices.

A rotation \code{R u q} has width $n_u$. For $q\ne0$, its axis must denote a
Hermitian unitary matrix: $A^\dagger=A$ and $A^\dagger A=I$.
Consequently $A^2=I$. This is stronger than merely being Hermitian
and is what makes the finite rotation formula possible.
At $q=0$, any well-formed axis is accepted: $e^{0A}=I$.
Its subterms are still checked, so a zero angle cannot conceal invalid syntax.
The adjoint $\dagger$ means conjugate transpose.

\begin{minipage}{\linewidth}
\begin{lstlisting}
data QOp
  = Id Nat                    -- identity on n qubits
  | Phase Rational            -- a scalar phase, on zero qubits
  | X | Y | Z | H | SX         -- fixed one-qubit gates
  | R QOp Rational            -- rotation about an operator axis
  | C QOp                     -- add a leading control qubit
  | Permute [Nat]             -- reorder the local wires
  | Tensor QOp QOp
  | DirectSum QOp QOp
  | Compose QOp QOp
  | Adjoint QOp
\end{lstlisting}
\end{minipage}

\paragraph{How code expressions are read.}
In QPP's infix notation, \code{u <.> t} is \code{Tensor u t},
\code{u <+> t} is \code{DirectSum u t}, and \code{u <> t} is
\code{Compose u t}.
For example, \code{Compose (C X) (Tensor H (Id 1))} applies a Hadamard
to the first qubit and then a controlled $X$.
Different terms can denote the same matrix: \code{Compose H H} and
\code{Id 1} are semantically, but not syntactically, equal.

\paragraph{Programs.}
QPP defines \code{Program = [Step]}, with \code{Unitary QOp},
\code{Initialize [Nat] [Bool]}, and \code{Measure [Nat]} as step constructors.
The exact prototype accepts all three constructors. Unitary steps have the
input width; measurement and reset retain unnormalized branches
(Section~\ref{sec:measurement}). All steps are validated before execution.
For $p=[\code{Unitary}\ u_1,\ldots,\code{Unitary}\ u_m]$ in execution order,
\[
\sem p\,v=\sem{u_m}\cdots\sem{u_1}\,v,\qquad \sem{[\,]}\,v=v.
\]

\section{One operator semantics for both evaluators}
\label{sec:matrix}
We first define meaning independently of storage and algorithms.
If $D_u=2^{n_u}$, each row of the following table specifies the denotation
of a constructor. Matrix multiplication implements composition;
$\otimes$ is the Kronecker product and $\diag(A,B)$ is a block diagonal
matrix. The rotation convention includes a half angle.
\begin{center}
\small
\begin{tabularx}{\linewidth}{@{}lY@{}}
\toprule
Term & Matrix denotation\\
\midrule
\code{Id n} & $I_{2^n}$.\\
\code{Phase q} & $[e^{i\pi q}]$.\\
\code{X, Y, Z, H, SX} & The matrices in \eqref{eq:primitives}.\\
\code{R u q} & $\exp(-i\pi q\,\sem u/2)$.\\
\code{C u} & $\diag(I_{D_u},\sem u)$.\\
\code{Permute ks} & $P_{f(j),j}=1$ and other entries zero; the bits of
 $f(j)$ are $b_{\code{ks}[0]},\ldots,b_{\code{ks}[n-1]}$ if $j$ has
 bits $b_0,\ldots,b_{n-1}$ and $n=|\code{ks}|$.\\
\code{Tensor u t} & $\sem u\otimes\sem t$.\\
\code{DirectSum u t} & $\diag(\sem u,\sem t)$, with $n_u=n_t$.\\
\code{Compose u t} & $\sem u\,\sem t$, with $n_u=n_t$.\\
\code{Adjoint u} & $\sem u^\dagger$.\\
\bottomrule
\end{tabularx}
\end{center}
\begin{equation}
\begin{gathered}
X:=\sem{\code X}=\begin{pmatrix}0&1\\1&0\end{pmatrix},\quad
Y:=\sem{\code Y}=\begin{pmatrix}0&-i\\i&0\end{pmatrix},\quad
Z:=\sem{\code Z}=\begin{pmatrix}1&0\\0&-1\end{pmatrix},\\
H:=\sem{\code H}=\frac1{\sqrt2}\begin{pmatrix}1&1\\1&-1\end{pmatrix},\quad
\mathrm{SX}:=\sem{\code{SX}}=\frac12
\begin{pmatrix}1+i&1-i\\1-i&1+i\end{pmatrix}.
\end{gathered}
\label{eq:primitives}
\end{equation}
The chosen square root of $X$ satisfies $\mathrm{SX}^2=X$.
It is unitary but not Hermitian, so \code{SX} is not a valid axis
at a nonzero angle. The term \code{H}, in contrast, is a valid axis:
rotations are not restricted to Pauli matrices.

\paragraph{Why global phase must be retained.}
For an isolated state, $v$ and $e^{i\alpha}v$ give the same measurement
probabilities. Inside an operator language, however, control converts
this phase into a relative phase:
\[
\sem{\code{C (Phase q)}}=
\begin{pmatrix}1&0\\0&e^{i\pi q}\end{pmatrix}.
\]
Thus equality includes phase: $\sem{\code{C (R (Id 0) 2)}}=Z$, not $I_2$.
Writing $S=\diag(1,i)$ and $T=\diag(1,e^{i\pi/4})$, the rotation convention gives
\[
\sem{\code{R Z (1/2)}}=e^{-i\pi/4}S,\qquad
\sem{\code{R Z (1/4)}}=e^{-i\pi/8}T.
\]
The terms \code{Phase (1/4) <.> R Z (1/2)} and
\code{Phase (1/8) <.> R Z (1/4)} denote $S$ and $T$ exactly.
These examples explain why a semantics suitable for compiler validation
must preserve scalar phases even on zero qubits.
\section{Cyclotomic arithmetic: what a coefficient contains}
\label{sec:scalar}
An exact coefficient is a complex number represented by rational
coordinates in a specified basis. We distinguish the number $z$, a
polynomial $p(x)$ representing it, and the tuple stored by the evaluator.
Here $x$ is a formal variable; $\zeta_N$ is an actual complex root of unity,
and substitution gives $z=p(\zeta_N)$. Several polynomials can represent
the same number. Reduction will select one of them uniquely.
Section~\ref{sec:implementation} adopts these coordinates as the baseline
storage format, with an explicit field context.

\subsection{A first example: the eighth roots of unity}
Let $\zeta_8=e^{2\pi i/8}=(1+i)/\sqrt2$. Its powers satisfy
$\zeta_8^2=i$ and $\zeta_8^4=-1$. The latter identity lets us replace
any higher power by a rational multiple of a lower one; for example
$\zeta_8^6=-\zeta_8^2$ and $\zeta_8^9=\zeta_8$.
Every element of the field $K_8=\Q(\zeta_8)$ has a unique expression
\[
z=z_0+z_1\zeta_8+z_2\zeta_8^2+z_3\zeta_8^3,\qquad z_j\in\Q.
\]
Here $\Q(\zeta_8)$ means the smallest field containing the rationals and
$\zeta_8$: it is closed under addition, subtraction, multiplication, and
division by nonzero elements. The four powers above form a basis over
$\Q$, so $z$ has coordinates $(z_0,z_1,z_2,z_3)$ in that basis.
The reason four coordinates suffice, even after division, is given below.

The Hadamard coefficient $h=1/\sqrt2$ is the real part of $\zeta_8$.
Since $\zeta_8^{-1}=-\zeta_8^3$, it has the exact representation
\[
h=\frac{\zeta_8+\zeta_8^{-1}}2
  =\frac{\zeta_8-\zeta_8^3}{2}
\quad\longleftrightarrow\quad (0,1/2,0,-1/2).
\]
Squaring this expression illustrates the arithmetic:
\[
h^2
=\frac{\zeta_8^2-2\zeta_8^4+\zeta_8^6}{4}
=\frac{\zeta_8^2+2-\zeta_8^2}{4}
=\frac12.
\]
We multiplied symbolic expressions and used $\zeta_8^4=-1$ to reduce
the powers. Thus the cancellation needed to verify $H^2=I$ uses rational
arithmetic; it never evaluates a decimal approximation to $\sqrt2$.

\subsection{Why finite coordinates exist and are unique}
For any positive integer $N$, fix $\zeta_N=e^{2\pi i/N}$ and write
$K_N=\Q(\zeta_N)$. This root is \emph{primitive}: its smallest positive
power equal to one is $N$. The primitive $N$th roots are precisely
$\zeta_N^k$ for $1\le k\le N$ with $\gcd(k,N)=1$.
The \emph{cyclotomic polynomial} $\Phi_N(x)$ has these roots, each once,
and leading coefficient one. For example, the primitive eighth roots
are the four odd powers of $\zeta_8$, and $\Phi_8(x)=x^4+1$.

The key theorem is that $\Phi_N$ has rational coefficients and is
irreducible over $\Q$: it has no factorization into two rational
polynomials of positive degree \cite{milne}. Its degree is
$d=\varphi(N)$, the number of integers in $\{1,\ldots,N\}$ coprime to $N$.
It is the polynomial of least positive degree over $\Q$ that vanishes
at $\zeta_N$.

To reduce a polynomial $p(x)$, divide it by $\Phi_N(x)$:
\[
p(x)=q(x)\Phi_N(x)+r(x),\qquad \deg r<d.
\]
Evaluating at $\zeta_N$ gives $p(\zeta_N)=r(\zeta_N)$ because
$\Phi_N(\zeta_N)=0$. The coefficients of $r$ are therefore sufficient.
They are also unique: two different remainders with the same value
would have a nonzero difference of degree below $d$ vanishing at
$\zeta_N$, contradicting minimality.

So far this establishes unique coordinates for polynomial expressions.
Within the fixed basis, equality is coordinatewise rational equality,
and zero means all coordinates are zero. We next show that division
also produces a polynomial expression, completing the field construction.

\subsection{Arithmetic is polynomial arithmetic followed by reduction}
Suppose $a=p(\zeta_N)$ and $b=q(\zeta_N)$, with $p,q$ already reduced.
For $a+b$, add their coefficient tuples. For $ab$, multiply $p$ and $q$,
then take the remainder on division by $\Phi_N$. The computation of
$h^2$ above is exactly this procedure for $\Phi_8=x^4+1$.

Conjugation uses the fact that a root of unity lies on the unit circle.
Since the coefficients of $p$ are rational and hence real,
\[
\overline{p(\zeta_N)}
=p(\overline{\zeta_N})
=p(\zeta_N^{-1})
=p(\zeta_N^{N-1}).
\]
Substitute $x^{N-1}$ for $x$ in $p$, then reduce as before.
For instance, $i=\zeta_8^2$ becomes $\zeta_8^6=-i$ under conjugation.

Division also stays in the same coordinate system. If $p\ne0$ and
$\deg p<d$, irreducibility implies $\gcd(p,\Phi_N)=1$.
The extended Euclidean algorithm on rational polynomials finds $s,t$ with
\[
s(x)p(x)+t(x)\Phi_N(x)=1.
\]
At $x=\zeta_N$, the second term vanishes, leaving
$s(\zeta_N)p(\zeta_N)=1$. Thus $s$, reduced modulo $\Phi_N$, represents
$a^{-1}$. Rational expressions therefore need no extra basis elements.
As a small example, $h^2=1/2$ gives $h^{-1}=2h$.

The polynomial values form a field containing $\Q$ and $\zeta_N$, and
every field containing these must contain those values. They are
therefore exactly $K_N$. In the standard quotient notation,
\begin{equation}
\Q[x]/(\Phi_N(x))\ \cong\ K_N,\qquad
[p]\longmapsto p(\zeta_N).
\label{eq:field}
\end{equation}
Here $\Q[x]$ denotes rational polynomials, and $[p]$ identifies all
polynomials differing from $p$ by a multiple of $\Phi_N$.
Its unique reduced representative supplies the stored coordinates.
A \emph{cyclotomic number} belongs to some $K_N$.

\subsection{Changing coordinates to combine fields}
Tuples from different fields cannot simply be added entry by entry:
their basis elements differ. For $a\in K_N$ and $b\in K_M$, choose
$L=\lcm(N,M)$. Our fixed choice of complex roots gives
\[
\zeta_N=\zeta_L^{L/N},\qquad \zeta_M=\zeta_L^{L/M}.
\]
Substitute these powers in the expressions for $a$ and $b$ and reduce
modulo $\Phi_L$. The resulting tuples share the basis
$1,\zeta_L,\ldots,\zeta_L^{\varphi(L)-1}$ and can be combined or compared.
These embeddings change the representation, not the complex number.

For example, a Hadamard coefficient in $K_8$ and a phase
$z=\zeta_6$ can both be represented in $K_{24}$. Put $\xi=\zeta_{24}$;
then $\zeta_8=\xi^3$, $\zeta_6=\xi^4$, and
$\Phi_{24}(x)=x^8-x^4+1$. Hence $\xi^9=\xi^5-\xi$, giving
\[
h=\frac{\xi^3-\xi^9}{2}
 =\frac{\xi+\xi^3-\xi^5}{2},\qquad z=\xi^4.
\]
In the eight-coordinate basis, $h$ is now
$(0,1/2,0,1/2,0,-1/2,0,0)$; the value is still $1/\sqrt2$.
Section~\ref{sec:example} uses precisely this pair of coefficients.

For a whole program we can choose a sufficient common field once.
Its index need not be minimal: $K_3=K_6$, since
$\zeta_6=1+\zeta_3$ and $\zeta_3=\zeta_6^2$.
The field degree counts rational coordinates, while their numerator
and denominator bit lengths determine the cost of rational arithmetic.
Both sizes matter; exactness alone gives no efficiency guarantee.

\section{Why every valid rational-angle term stays cyclotomic}
\label{sec:closure}
Rational-angle phases are roots of unity; rotations need a further argument.
For a nonzero rotation angle, the unitary axis must also be Hermitian.
For $A=\sem u$ this implies $A^2=I_{2^{n_u}}$. Therefore every even power
of $A$ is identity and every odd power is $A$. With $\theta=\pi q/2$,
the matrix exponential separates into two scalar series:
\begin{equation}
\begin{aligned}
e^{-i\theta A}
&=\left(\sum_{k\ge0}\frac{(-1)^k\theta^{2k}}{(2k)!}\right)I
 -i\left(\sum_{k\ge0}\frac{(-1)^k\theta^{2k+1}}{(2k+1)!}\right)A\\
&=\cos\theta\,I-i\sin\theta\,A.
\end{aligned}
\label{eq:exp}
\end{equation}
For a reduced rational $q=a/b$, $b>0$, set
$w=e^{i\theta}=\zeta_{4b}^{\,a}$. Euler's identities now give the finite
formula used by both evaluators:
\begin{equation}
\begin{gathered}
\sem{\code{R}\ u\ q}
 =c_q\,\sem{\code{Id}\ n_u}+d_q\,\sem u,\\
c_q=\frac{w+w^{-1}}2,\qquad d_q=\frac{w^{-1}-w}2.
\end{gathered}
\label{eq:rotation}
\end{equation}
The series proves the identity; the evaluator never computes the series,
trigonometric functions, or a numerical matrix exponential.
At $q=0$, $c_q=1$ and $d_q=0$, so the same formula holds without
assuming a Hermitian axis; evaluation skips the axis action.
For example, $q=1/3$ gives $c_q=\sqrt3/2$ and $d_q=-i/2$.
The formula applies to any valid composite axis, including \code{H}
and \code{C X}, not only to single-qubit Pauli axes.

\paragraph{Choosing a sufficient field.}
A phase with parameter $a/b$ is $\zeta_{2b}^{\,a}$.
The fixed gates need only $i$ and $1/\sqrt2$, both in $K_8$.
Write $s\preceq t$ when $s$ is a syntactic subterm of $t$, including $t$.
Collect all angle denominators, including those inside rotation axes:
\begin{equation}
N=\lcm\!\left(\{8\}\cup\{2b:\code{Phase}\ (a/b)\preceq t\}
                    \cup\{4b:\code{R}\ u\ (a/b)\preceq t\}\right).
\label{eq:conductor}
\end{equation}
For instance, occurrences of \code{Phase (1/3)} and \code{R Z (1/5)}
give the sufficient index $\lcm(8,6,20)=120$, of degree $\varphi(120)=32$.
This index need not be the minimal conductor.
After validating all subterms, the prototype replaces zero-angle rotations
by identities before applying this rule; unused axes contribute no denominators.

\begin{proposition}
Every valid $n$-qubit term $t$ denotes a unitary matrix with entries in
$D_N=\Z[1/2,\zeta_N]\subset K_N$, for $N$ from \eqref{eq:conductor}.
\end{proposition}
\begin{proof}
The ring $D_N$ consists of finite sums of powers of $\zeta_N$ with
dyadic rational coefficients (integers divided by powers of two).
The fixed gates and phases lie in this ring; \eqref{eq:rotation} uses
only its addition, multiplication, and halving operations.
Composition, tensor and direct sums, control, permutations, and adjoints
preserve the ring, since conjugation sends $\zeta_N$ to $\zeta_N^{-1}$.
Structural induction gives closure. The same induction proves unitarity:
each primitive is unitary, and exponentiating a Hermitian axis with a
purely imaginary parameter is unitary. A zero-angle rotation is identity.
\end{proof}
If an input vector has entries in $K_M$, application stays in
$K_{\lcm(N,M)}$. General field inverses are unnecessary for unitary
evaluation, but will be needed for MPS rank factorization.
Exact synthesis asks the converse kind of question: Giles and Selinger
characterize the matrices synthesizable over Clifford+$T$ with ancillas
using $\Z[1/\sqrt2,i]$ \cite{giles}. Here a given expression determines
its field; arbitrary rational angles need not stay in that fixed ring.

\section{Worked evaluation: a phase and a non-Pauli rotation}
\label{sec:example}
Consider a Bell state followed by a phase gate on its second qubit.
The following Haskell definitions build the operator expression:
\begin{lstlisting}
phaseGate q = C (Phase q)
bell = Compose (C X) (Tensor H (Id 1))
example = Compose (Tensor (Id 1) (phaseGate (1/3))) bell
\end{lstlisting}
\code{phaseGate q} controls a zero-qubit scalar. Let $h=1/\sqrt2$,
$z=e^{i\pi/3}=\zeta_6$, and
$B=\sem{\code{bell}}$. In the basis $(00,01,10,11)$, the matrix evaluator
constructs
\[
B=h\begin{pmatrix}
1&0&1&0\\0&1&0&1\\0&1&0&-1\\1&0&-1&0
\end{pmatrix},\qquad
U:=\sem{\code{example}}
=h\begin{pmatrix}
1&0&1&0\\0&z&0&z\\0&1&0&-1\\z&0&-z&0
\end{pmatrix}.
\]
The second matrix is obtained by multiplying $B$ on the left by
$I_2\otimes\diag(1,z)=\diag(1,z,1,z)$.
The direct evaluator instead starts with a length-four vector and performs
three local actions:
\begin{equation}
\ket{00}
\ \longmapsto\ h(\ket{00}+\ket{10})
\ \longmapsto\ h(\ket{00}+\ket{11})
\ \longmapsto\ h(\ket{00}+z\ket{11}).
\label{eq:bell}
\end{equation}
The last vector is exactly the first column of $U$. No $4\times4$ matrix
is needed on this execution path.

\paragraph{The actual coefficient field.}
The expression contains Hadamards and \code{Phase (1/3)}, so the field
rule gives $N=\lcm(8,6)=24$, with $\varphi(24)=8$.
Both backends can use the basis modulo $\Phi_{24}(x)=x^8-x^4+1$.
In that common field,
\[
z=\zeta_{24}^{\,4},\qquad
h=\frac{\zeta_{24}^{\,3}+\zeta_{24}^{-3}}2,\qquad
\|U\ket{00}\|^2=h^2(1+\bar z z)=1.
\]

\paragraph{Observing the relative phase.}
Applying the inverse Bell circuit gives
\[
B^\dagger U\ket{00}
 =\frac{1+z}{2}\ket{00}+\frac{1-z}{2}\ket{10}.
\]
Since $z+\bar z=1$, the squared amplitudes are
$|1+z|^2/4=3/4$ and $|1-z|^2/4=1/4$.
Discarding the phase gate would give $\ket{00}$. These exact probabilities
require neither random sampling nor state collapse.

\paragraph{A rotation about the Hadamard axis.}
Let $r=\code{R H (1/3)}$. Since $H^2=I$, \eqref{eq:rotation} gives
$w=\zeta_{12}$, $c=\sqrt3/2$, and $d=-i/2$.
The common field is $K_{24}$: in the notation of Section~\ref{sec:scalar},
$w=\xi^2$ and $i=\xi^6$. The matrix evaluator constructs
\[
\sem r=cI_2+dH=
\begin{pmatrix}c+dh&dh\\dh&c-dh\end{pmatrix}.
\]
The direct evaluator computes $H\ket0=h(\ket0+\ket1)$ and combines it
with the original input:
\[
\sem r\,\ket0
=c\ket0+d\,h(\ket0+\ket1)
=\left(\frac{\sqrt3}{2}-\frac{i}{2\sqrt2}\right)\ket0
 -\frac{i}{2\sqrt2}\ket1.
\]
The squared magnitudes of these first-column entries are $7/8$ and $1/8$,
so the norm is one. This exercises the half angle, sign, and recursive axis action.

\section{Shared implementation foundations}
\label{sec:implementation}
With \code{St = QPP.Exact.State} and
\code{SV = QPP.Semantics.Exact.Statevector}, evaluate the preceding example:
\begin{lstlisting}[label=lst:public-example]
evaluateExample = do
  let input = St.ket [0, 0]
      output = SV.apply (SV.evalOp example) input
  norm2 <- St.normSquared defaultLimits output
  pure (St.amplitudes output, norm2)
\end{lstlisting}
The result is \code{Right} $([h,0,0,hz],1)$ in \code{Cyclo} values.
Both backends share \code{StateT}. Ordinary \code{ket}, \code{evalOp},
\code{apply}, and \code{evalProg} match the numerical signatures, use
\code{defaultLimits}, and raise runtime errors on failure. Their
\code{WithLimits} variants take explicit limits and return \code{Result}.
\code{evalProg program input rng} returns \code{(output, [], rng)} for a
unitary program. Run \code{cabal run ExactEvaluation} for an example.

\subsection{Representation and field ownership}
An immutable \code{Field} records $N$, $\Phi_N$, and the distinguished
embedding. A \code{Coeff} stores $d=\varphi(N)$ reduced rational coordinates,
padded with zeros; \code{Cyclo} pairs them with their field.
Arrays share that context:
\begin{lstlisting}
data Buffer = Buffer !Field !(V.Vector Coeff)
data StateT = StateT !Int !Buffer
\end{lstlisting}
\code{V.Vector} is a boxed array. A \code{StateT n} buffer has $2^n$
coefficients in its field's basis: $D$ amplitudes occupy $Dd$ rational slots.
Matrices use the same buffer with a different shape. Internal constructors
and checked imports enforce these invariants.

Public operations promote to the least-common-multiple field; kernels
keep it. In listings, \code{C} names \code{QPP.Exact.Checked}, \code{S}
names scalar arithmetic, and \code{Result a} is \code{Either ExactError a}.
Listings reproduce constructor recursions, omitting imports and deriving
clauses. Adaptations and omitted helpers are identified where used.

\subsection{Validation and resource limits}
\code{C.checkOp} checks widths, permutations, and Hermiticity of nonzero-angle
rotation axes. Its symbolic rules cover primitives and supported composite
forms; rejection does not prove non-Hermiticity. Thus \code{R H q} is
accepted; \code{R SX q} requires $q=0$. An exact matrix test
$A=A^\dagger$ is a possible future fallback.

The private \code{CheckedOp} tree caches widths, dimensions, and field indices.
Write $\check u$ for the successful validation of source term $u$;
$\sem{\check u}=\sem u$. Validation precedes pruning zero-angle axes.
\code{C.view} exposes the checked constructors: \code{Id} gets its width
from the node, and \code{X,Y,Z,H,SX} become \code{Primitive g}.
\code{C.toQOp} recovers the validated, pruned syntax for the dense evaluator.

Limits bound root index, degree, coordinate slots, and integer sizes.
Preflight uses unbounded dimension arithmetic and checks $N$ before
factorization. Workspace includes constants, retained arrays, promotion,
and temporaries. Conservative checks may reject later cancellation;
they bound neither wall time nor heap bytes. Errors distinguish invalid
input, width mismatch, division by zero, and exhaustion, with term and
zero-based program-step context. Even equality can fail; exhaustion
never means inequality or approximation.

\subsection{Linear algebra and the normalization boundary}
\code{HilbertSpace} supplies scaling, addition, subtraction, and \code{inner};
squared norm is \code{inner v v}. Numerical Matrix uses \code{CMat} for
both states and operators; Statevector supplies the instance for states.
The class additionally requires
\begin{lstlisting}
Scalar v ~ Complex (Realnum v), Floating (Realnum v)
norm      :: v -> Realnum v
normalize :: v -> v
\end{lstlisting}
For $a,b\in K$ and $v,w\in K^D$, the field-compatible operations are
\[
av+bw\in K^D,\qquad
\langle v,w\rangle=\sum_j\overline{v_j}w_j\in K,\qquad
\langle v,v\rangle\in K\cap\mathbb R_{\ge0}.
\]
The prototype implements these directly in \code{QPP.Exact.State}:
\begin{lstlisting}
scale       :: Limits -> Cyclo -> StateT -> Result StateT
add, sub    :: Limits -> StateT -> StateT -> Result StateT
inner       :: Limits -> StateT -> StateT -> Result Cyclo
normSquared :: Limits -> StateT -> Result Cyclo
\end{lstlisting}
The dense backend has corresponding operator operations, with
$\langle A,B\rangle=\tr(A^\dagger B)$. Linear combinations need not be
unitary, so only construction by \code{evalOp} guarantees unitarity.
A future class should separate these capabilities from normalization,
using an associated scalar type. The prototype claims no \code{HilbertSpace}
instance. Cyclotomic scalars are already complex; \code{Complex Cyclo}
and \code{Data.Complex}'s \code{RealFloat} constraint are inappropriate
\cite{base}.

Arbitrary normalization is an algebraic obstruction, not just an interface
problem: $(1,2+\sqrt2)$ has squared norm $7+4\sqrt2$, whose square root
is noncyclotomic. Unitary evaluation preserves the input norm and needs
no normalization.

\paragraph{Optional proof of the normalization boundary (Galois theory).}
Every subfield of a cyclotomic field is Galois over $\Q$: its Galois group
is a quotient of an abelian group, whose subgroups are all normal
\cite{milne}. Put $F=\Q(\sqrt2)$, $a=7+4\sqrt2$, and $a'=7-4\sqrt2$.
The field norm gives $aa'=17$, so neither $a$ nor $a'$ is a square in $F$:
the norm of a square would be a rational square.
For $\sqrt{a'}$ to lie in $F(\sqrt a)$, the ratio
$a'/a=(a')^2/17$ would have to be a square in $F$, but $\sqrt{17}\notin F$.
Thus $\Q(\sqrt a)=F(\sqrt a)$ lacks a conjugate of $\sqrt a$ and is not
Galois. Neither $\sqrt a$ nor its reciprocal can be cyclotomic.
\label{sec:normalization}

Measurement avoids this obstruction by retaining projected vectors and using
ratios of squared norms. Section~\ref{sec:measurement} extends the following
unitary evaluators to measurement and reset without changing scalar fields.

\Needspace{18\baselineskip}
\section{Constructing an exact dense matrix}
\label{sec:dense}
\code{QPP.Semantics.Exact.Matrix} constructs $\sem u$ explicitly.
Its internal \code{OpT n buffer} stores $2^{2n}$ entries in row-major order.
\code{evalOpWithLimits} calls \code{C.checkOp}, selects \code{C.fieldIndex checked}
using \code{S.makeField}, checks \code{matrixBudget}, then invokes
\code{evaluate limits field checked}.

\Needspace{9\baselineskip}
\subsection{The recursive evaluator}
This excerpt preserves every constructor case. \code{M} names the matrix
kernels; the omitted \code{pack} attaches width and field to their array.
Omitted local helpers lift Kronecker product ($\otimes$), multiplication
($\circ$), direct sum (\code{<+>}), addition (\code{.+}), scaling (\code{.*}),
and adjoint (\code{adj}) through \code{Result}, propagating failures.
Their mathematical operations are specified below.
\begin{lstlisting}[label=lst:matrix-evaluate]
evaluate :: Limits -> Field -> C.CheckedOp -> Result OpT
evaluate limits field checked =
  pack (C.qubits checked) field <$> eval (C.toQOp checked)
  where
    eval op = case op of
      Id n    -> ident (2^n)
      Phase q -> phase q
      X  -> gate C.X
      Y  -> gate C.Y
      Z  -> gate C.Z
      H  -> gate C.H
      SX -> gate C.SX
      R axis q -> do
        (c,d) <- S.rotationC limits field q
        if S.isZeroC d
          then c .* ident (2 ^ op_qubits axis)
          else do
            a <- eval axis
            (c .* ident (M.dimension a)) .+ (d .* pure a)
      Permute ks -> pure (M.permutation field ks)
      C a -> do
        mat <- eval a
        ident (M.dimension mat) <+> pure mat
      Tensor a b    -> eval a ⊗ eval b
      Compose a b   -> eval a ∘ eval b
      DirectSum a b -> eval a <+> eval b
      Adjoint a     -> adj (eval a)
    ident dim = pure (M.identity field dim)
    phase q = M.scalar <$> S.phaseC limits field q
    gate g = M.Matrix 2 <$> primitive limits field g
\end{lstlisting}
\code{primitive} supplies \eqref{eq:primitives}; \code{rotationC} implements
\eqref{eq:rotation}. When $d=0$, the axis
need not be evaluated. A phase subterm is $1\times1$; control or direct
sum incorporates it as a block in a larger matrix. In particular,
\code{C (Phase q)} gives $\diag(1,e^{i\pi q})$, of dimension two.

\subsection{Application, correctness, and cost}
For compatible matrices and a column vector, the kernels compute
\[
(AB)_{ij}=\sum_k A_{ik}B_{kj},\qquad
(Av)_i=\sum_j A_{ij}v_j,\qquad
(A\otimes B)_{ir+k,\,js+l}=A_{ij}B_{kl},
\]
where $B$ has shape $r\times s$. Conjugate transpose exchanges indices
and conjugates coefficients. Accumulation starts at exact zero, and
entries are skipped only after exact zero tests. \code{apply} checks
widths, promotes matrix and input to a common field, and computes $Av$.
\code{evalProg} selects one field for the entire program; the shared runner
applies unitaries, measurements, and resets in execution order.

Structural induction shows that successful evaluation of $\check u$ stores
the entries of $\sem u$, providing an oracle for the other backends.
For $D=2^n$, the output occupies $O(D^2)$ field elements;
multiplication costs $O(D^3)$ field operations and application $O(D^2)$.
Entrywise equality includes global phase and reduces to rational equality
after promotion. These operation counts are not bit complexities.

\Needspace{18\baselineskip}
\section{Direct exact statevector evaluation}
\label{sec:sv}
For a valid source term $u$, put $D_u=2^{n_u}$ and choose a common cyclotomic
field $K$ containing its coefficients and those of the input.
The direct evaluator computes $\sem u\,v$ for $v\in K^{D_u}$ without
constructing $\sem u$. Vectors may be unnormalized or entangled, in the
computational-basis order of Section~\ref{sec:language}.

\Needspace{9\baselineskip}
\subsection{Deriving the direct actions}
\paragraph{Identity, phase, and one-qubit gates.}
With $z_q=e^{i\pi q}$, the elementary actions are
\[
\sem{\code{Id}\ n}\,v=v,\qquad
\sem{\code{Phase}\ q}\,(a)=(z_q a),\qquad
G\binom{a}{b}=
\binom{G_{00}a+G_{01}b}{G_{10}a+G_{11}b}.
\]
Here $G$ is the matrix of a primitive or its prepared adjoint.

\paragraph{Rotation and composition.}
Distributing matrix application over the finite rotation formula gives
\begin{equation}
\sem{\code{R}\ u\ q}\,v
=(c_qI+d_q\sem u)v
=c_qv+d_q(\sem u\,v).
\label{eq:svrotation}
\end{equation}
First compute $w=\sem u\,v$, retaining $v$, then combine $c_qv_j+d_qw_j$
entrywise. This handles \code{H}, \code{C X}, and composite axes alike.
Composition applies the right child first:
\[
\sem{\code{Compose}\ u\ t}\,v=\sem u\,(\sem t\,v).
\]

\paragraph{Direct sum and control.}
Write $v=v_0\mathbin{\|}v_1$ for concatenation of equal-length vectors.
The leading qubit selects the two contiguous blocks. Block multiplication gives
\begin{align}
\sem{\code{DirectSum}\ u\ t}\,(v_0\mathbin{\|}v_1)
 &=(\sem u\,v_0)\mathbin{\|}(\sem t\,v_1),\\
\sem{\code{C}\ u}\,(v_0\mathbin{\|}v_1)
 &=v_0\mathbin{\|}(\sem u\,v_1).
\end{align}
The second rule is the first with identity on the zero block.
In particular, $\sem{\code{C (Phase q)}}(a,b)^T=(a,z_qb)^T$.

These blocks must retain their norms. The phased Bell state of
\eqref{eq:bell} has branches $(h,0)^T$ and $(0,hz)^T$, each with squared
norm $1/2$. Normalizing them separately would change the state and break
linearity. No recursive rule normalizes a vector.

\paragraph{Tensor application and memory layout.}
\label{sec:layout}
In general an entangled input cannot be written $v=x\otimes y$.
For \code{Tensor u t}, put $A=\sem u$ and $B=\sem t$.
View the full vector as the $D_u\times D_t$ array
$V_{ij}=v_{iD_t+j}$. The Kronecker definition gives
\begin{equation}
\bigl(\sem{\code{Tensor}\ u\ t}\,v\bigr)_{iD_t+j}
 =\sum_{k=0}^{D_u-1}\sum_{\ell=0}^{D_t-1}A_{ik}B_{j\ell}V_{k\ell}.
\label{eq:tensoraction}
\end{equation}
Separate the two sums:
\[
W_{ij}=\sum_{\ell=0}^{D_t-1}B_{j\ell}V_{i\ell},
\qquad
V'_{ij}=\sum_{k=0}^{D_u-1}A_{ik}W_{kj}.
\]
The first stage applies $B$ to each contiguous row, passed as a column
vector. The second applies $A$ to each gathered column. In matrix notation,
\begin{equation}
W=VB^T,\qquad V'=AW=AVB^T,\qquad
w_{iD_t+j}=V'_{ij}.
\label{eq:svlayout}
\end{equation}
This is an ordinary transpose: the index $j$ selects the output of $B$;
no coefficient is conjugated. These stages use the child transformations
without constructing $A$, $B$, or $A\otimes B$. A zero-qubit factor has
dimension one, so the same rule includes scalar tensor factors.

For example, let $T=\diag(1,\zeta_8)$ and $v=(a,b,c,d)^T$.
The row step for $H\otimes T$ gives $(a,\zeta_8b)$ and
$(c,\zeta_8d)$; the column step gives
\[
(H\otimes T)v=\frac1{\sqrt2}
(a+c,\ \zeta_8(b+d),\ a-c,\ \zeta_8(b-d))^T.
\]

\paragraph{Permutations.}
Let $f$ be the basis-index permutation specified by \code{ks}.
The matrix rule sends $\ket j$ to $\ket{f(j)}$. Hence
\[
\bigl(\sem{\code{Permute ks}}\,v\bigr)_j=v_{f^{-1}(j)}.
\]
For example, \code{Permute [1,2,0]} maps bits $abc$ to $bca$:
output index $1$ receives input index $4$.
For forward execution, the prepared \code{Permute} stores the inverse
permutation of the $n$ \emph{wire positions}. \code{Vector.permute} computes the induced
basis-index gather $f^{-1}(j)$ on demand; no $2^n$-entry index table is stored.

\paragraph{Preparing adjoints.}
An opaque vector-transforming function does not expose its adjoint.
We retain the expression structure during preparation and use
\[
(AB)^\dagger=B^\dagger A^\dagger,\qquad
(A\otimes B)^\dagger=A^\dagger\otimes B^\dagger,\qquad
(A\oplus B)^\dagger=A^\dagger\oplus B^\dagger.
\]
Control follows from $I\oplus A$. At the leaves, phases use
$\overline{z_q}=z_{-q}$, gate tables use
$(G^\dagger)_{ij}=\overline{G_{ji}}$, and permutations invert $f$.
For a nonzero-angle rotation, $c_q$ is real, $d_q$ is purely imaginary,
and the checked axis $A$ is Hermitian, so
\begin{equation}
\sem{\code{R}\ u\ q}^{\dagger}
=(c_qI+d_qA)^\dagger=c_qI-d_qA
=\sem{\code{R}\ u\ (-q)}.
\label{eq:svadjoint}
\end{equation}
These laws let preparation push adjoints to leaves. For rotations, the
axis remains unchanged. Zero-angle rotations were already validated
and replaced by identity.

\subsection{The prepared representation}
\label{sec:plans}
The prepared tree specializes \code{QOp} to one field: angles and primitive
symbols become coefficients, and adjoints are pushed to leaves.
It supports coefficient reuse and promotion; validated syntax with a
coefficient cache is an alternative.
\code{QPP.Exact.Internal.Plan} defines the datatype and \code{compile};
\code{QPP.Semantics.Exact.Statevector} defines \code{runPlan}.
\begin{lstlisting}[label=lst:plan-type]
data Plan a = Plan !Int (Action a)
data Action a
  = Id | Scalar a | Gate (V.Vector a)
  | R (Plan a) a a | C (Plan a) | Permute [Int]
  | Tensor (Plan a) (Plan a)
  | DirectSum (Plan a) (Plan a)
  | Compose (Plan a) (Plan a)
\end{lstlisting}
\code{Plan dim action} records dimension: composition preserves it,
tensor multiplies it, and direct sum adds equal dimensions. \code{Coeff}
uses \code{OpT}'s field. Derived \code{Foldable} and \code{Traversable}
instances visit constants for validation and embedding.

\subsection{Compiling a reusable operator}
With \code{P} naming the plan module, \code{OpT}, shown positionally,
stores width, conservative
counts of constants and live workspace (\code{Requirements}), field, and
plan. The listing replaces record construction by positional construction;
the preparation steps are unchanged:
\begin{lstlisting}[label=lst:plan-evalop,breaklines=true]
data OpT = OpT Int P.Requirements Field (P.Plan Coeff)
evalOpWithLimits :: Limits -> QOp -> Result OpT
evalOpWithLimits limits op = do
  checked <- C.checkOp limits op
  field <- S.makeField limits (C.fieldIndex checked)
  let required = P.requirements checked
  P.budget limits field (C.dimension checked) required
  plan <- P.compile limits field checked
  pure (OpT (C.qubits checked) required field plan)
\end{lstlisting}
\code{OpT} retains the plan, not a second copy of the checked tree.

\paragraph{The compiler.}
\code{compile} constructs the plan recursively. \code{Forward} prepares
$\sem u$ and \code{Backward} prepares $\sem u^\dagger$.
The helper \code{directedAngle} keeps the angle in \code{Forward} and
negates it in \code{Backward}; \code{opposite} exchanges the directions.
\code{phaseC}, \code{rotationC}, and \code{primitive} construct the exact
coefficients of Sections~\ref{sec:matrix} and~\ref{sec:closure}.
\Needspace{32\baselineskip}
\begin{lstlisting}[label=lst:plan-compile]
data Direction = Forward | Backward
compile :: Limits -> Field -> C.CheckedOp -> Result (Plan Coeff)
compile limits field = eval Forward
  where
    eval direction op = Plan (C.dimension op) <$> case C.view op of
      C.Id -> pure Id
      C.Phase q -> Scalar <$> S.phaseC limits field (directedAngle direction q)
      C.Primitive g -> Gate <$> gate direction g
      C.R axis q -> do
        (c,d) <- S.rotationC limits field (directedAngle direction q)
        if S.isZeroC d then pure (Scalar c)
        else do
          a <- eval Forward axis
          pure (R a c d)
      C.C a -> C <$> eval direction a
      C.DirectSum a b -> DirectSum <$> eval direction a <*> eval direction b
      C.Tensor a b -> Tensor <$> eval direction a <*> eval direction b
      C.Compose a b -> case direction of
        Forward -> Compose <$> eval Forward a <*> eval Forward b
        Backward -> Compose <$> eval Backward b <*> eval Backward a
      C.Adjoint a -> do
        Plan _ action <- eval (opposite direction) a
        pure action
      C.Permute ks -> pure $ Permute $ case direction of
        Forward -> inversePermutation ks
        Backward -> ks
    gate Forward g = primitive limits field g
    gate Backward g = do
      entries <- primitive limits field g
      V.generateM 4 $ \index ->
        let (row,column) = index `quotRem` 2
        in S.conjugateC limits field (entries V.! (2*column + row))
\end{lstlisting}
\begin{samepage}
The outer \code{Plan} preserves dimension, including for scalar rotations.
Hermitian axes compile \code{Forward} in either direction: only $d$
changes sign in \eqref{eq:svadjoint}. Adjoint reverses composition and
leaves no node; permutations become gather maps.
\par
\end{samepage}

\subsection{Executing and applying the plan}
\code{Vector.Transform} is \code{V.Vector Coeff -> Result (V.Vector Coeff)}.
\code{runPlan} returns this function; its helpers implement the actions
above. The listing expands the implementation's composition and direct-sum notation:
\begin{lstlisting}[label=lst:plan-run]
runPlan :: Limits -> Field -> P.Plan Coeff -> Vector.Transform
runPlan limits field = eval
  where
    eval (P.Plan _ action) = case action of
      P.Id -> pure
      P.Scalar c -> Vector.scale limits field c
      P.Gate g -> Vector.applyGate limits field g
      P.R axis c d -> \v -> do
        av <- eval axis v
        Vector.linearCombination limits field c v d av
      P.C a -> Vector.directSum pure (eval a)
      P.DirectSum a b -> Vector.directSum (eval a) (eval b)
      P.Compose a b -> \v -> eval b v >>= eval a
      P.Tensor a b -> Vector.tensor (P.dimension a) (P.dimension b)
                                    (eval a) (eval b)
      P.Permute gather -> pure . Vector.permute gather
\end{lstlisting}
\code{P.dimension (P.Plan dim action) = dim} supplies tensor shapes.
Composition passes the right child's output to the left. Execution
computes new amplitudes, but does not recompute gate coefficients,
construct composite matrices, or select a field.

\Needspace{15\baselineskip}
\paragraph{A construction and execution trace.}
For the worked \code{R H (1/3)} example, \code{compile} selects the
\code{C.R} case and recursively prepares \code{H}. In the field $K_{24}$,
let \code{hs} store $(h,h,h,-h)$, $h=1/\sqrt2$, and let
\code{c}, \code{d} store the coordinates of $c=\sqrt3/2$, $d=-i/2$.
Compilation constructs these nodes:
\begin{lstlisting}
pH = Plan 2 (Gate hs)
p  = Plan 2 (R pH c d)
\end{lstlisting}
Interpreting their stored coordinates as field elements gives
\[
\operatorname{runPlan}_{K_{24}}(\code p,(1,0)^T)
=c(1,0)^T+d\operatorname{runPlan}_{K_{24}}(\code{pH},(1,0)^T)
=(c+dh,dh)^T.
\]
\code{Gate} returns $(h,h)^T$; \code{R} combines it with the retained input.
The six constants are reused across applications, including tensor rows
and columns. Identical subtrees are not deduplicated.

\paragraph{Application, promotion, and programs.}
\code{apply} validates widths, state, and constants under its current limits.
It promotes operator and input from $K_N,K_M$ to $K_{\lcm(N,M)}$, rechecking
workspace, then calls the shared-field interpreter:
\begin{lstlisting}[label=lst:plan-apply]
applyIn :: Limits -> Field -> P.Plan Coeff -> StateT -> Result StateT
applyIn limits field plan (St.StateT n (Buffer _ input)) = do
  output <- runPlan limits field plan input
  pure (St.StateT n (Buffer field output))
\end{lstlisting}
\code{St} names the internal state module. Same-field plans are reused;
promotion is temporary, so repeated larger-field applications repeat embedding.
\code{evalProg} instead preflights all steps in one field, compiling each
unitary as it is reached. Measurement and reset use the shared interpreter.
An empty program preserves its validated input and RNG.

\subsection{Correctness and cost}
Let $\operatorname{prep}_K(\delta,\check u)$ denote the compiler's local
\code{eval direction op}. Its two-direction induction invariant is
\begin{equation}
\begin{aligned}
\operatorname{runPlan}_K(\operatorname{prep}_K(\code{Forward},\check u),v)
 &=\sem u\,v,\\
\operatorname{runPlan}_K(\operatorname{prep}_K(\code{Backward},\check u),v)
 &=\sem u^\dagger v.
\end{aligned}
\label{eq:planmeaning}
\end{equation}
\code{compile} selects \code{Forward}, so its execution gives $\sem u\,v$.
Here $K$ contains gate and input coefficients; limits and errors are
suppressed to describe successful execution.
\begin{proposition}
The ideal direct evaluator returns $w=\sem u\,v$ for every accepted term $u$
and compatible vector $v$. If $v$ has cyclotomic entries, so does $w$.
Moreover $w^\dagger w=v^\dagger v$.
\end{proposition}
\begin{proof}
Induct over the checked expression in both directions. Each interpreter
clause realizes its vector recurrence. Backward composition reverses its
children; backward rotation changes the angle but prepares the Hermitian
axis forward. The adjoint laws establish \eqref{eq:planmeaning}.
Closure follows from field arithmetic,
and $w^\dagger w=v^\dagger\sem u^\dagger\sem u\,v=v^\dagger v$.
\end{proof}
\paragraph{Promotion preserves execution.}
Let $\iota:K\hookrightarrow L$ be the embedding of Section~\ref{sec:scalar},
applied to vector entries and plan constants, preserving shape. Then
\begin{equation}
\iota\bigl(\operatorname{runPlan}_K(p,v)\bigr)
=\operatorname{runPlan}_L(\iota(p),\iota(v)).
\label{eq:promotion}
\end{equation}
For compatible $p,v$ over $K$, induction proves \eqref{eq:promotion} because
embeddings preserve arithmetic and commute with indexing and concatenation.
They preserve the selected roots, conjugation, and zero tests, so promotion
of a compiled plan agrees with compilation in $L$. Hence \code{apply}'s
promotion from operator/input fields $K_N,K_M$ to $K_{\lcm(N,M)}$ preserves
the result. Successive application proves the unitary program semantics;
Section~\ref{sec:measurement} establishes measurement and reset correctness.

\paragraph{Execution cost.}
After preparation, $G$ gates of bounded arity on $n$ qubits cost
$O(G2^n)$ field operations. In a sequential gate circuit, working vectors
need $O(2^n)$ live field elements, excluding stored plan constants.
The prototype allocates immutable arrays; this is a live-storage bound,
not a claim of buffer reuse. For an expanded expression tree with $s$
nodes, execution costs $O(s2^n)$ field operations, with indexing overhead.
Nested rotations may retain $O(h2^n)$ live elements at depth $h$.
Field construction, compilation, validation, and promotion are additional
costs; field degree and integer sizes determine bit costs.

\section{Measurement and reset without normalization}
\label{sec:measurement}
\paragraph{Stored vector and physical state.}
For nonzero $v\in K^{2^n}$, the physical state is $\rho_v=vv^\dagger/(v^\dagger v)$;
the implementation stores $v$. Let $P_b^{(k)}$ project onto bit $b$ on wire $k$.
The Born probabilities and stored branch are
\begin{equation}
 w_b=\|P_b^{(k)}v\|^2,\qquad
 p_b=\frac{w_b}{\langle v,v\rangle}=\frac{w_b}{w_0+w_1},\qquad v_b=P_b^{(k)}v.
 \label{eq:measurement}
\end{equation}
These quantities remain cyclotomic; the zero vector cannot be measured.
For a selected nonzero branch,
$\rho_{v_b}=P_b^{(k)}\rho_vP_b^{(k)}/p_b$, exactly the normalized conditional
density matrix. Later operations therefore have the usual physical meaning
without square roots. Replacing $v$ by $cv$, $c\ne0$, preserves every $p_b$.
However, \code{amplitudes}, \code{inner},
\code{normSquared}, and \code{equal} still describe the stored vector;
they do not implicitly normalize or compare rays. The checked function
\code{probability limits state wire outcome} computes $p_b$.

\paragraph{Sampling and the program interface.}
With outcome list $o$ and RNG stream $r:s$, the single-wire transition is
\[
 \sem{\code{Measure [k]}}(v,o,r:s)=(v_b,b:o,s),\qquad
 b=\begin{cases}0&r(w_0+w_1)<w_0,\\1&\text{otherwise}.\end{cases}
\]
Each call consumes one finite \code{Double} in $[0,1)$, even when an outcome
is certain. It treats $r$ as its exact dyadic rational; equality selects one.
This retains QPP's finite-precision RNG convention, while the threshold
comparison is exact. \code{compareReal} checks reality by conjugation and
equality algebraically. For a nonzero real value
$\sum_j c_j\zeta_N^j$, certified dyadic bounds enclose
$\sum_j c_j\cos(2\pi j/N)$ until its sign is determined.
Outward-rounded arithmetic bounds the alternating
\href{https://dlmf.nist.gov/4.19#E2}{cosine} and
\href{https://dlmf.nist.gov/4.24#E3}{arctangent} series, and
$\pi=16\arctan(1/5)-4\arctan(1/239)$. The sampler uses the same sign kernel
on the known-real difference $w_0-r(w_0+w_1)$.
Exhausting the bit limit reports an error, never a guessed outcome.

Program execution is shared in \code{Internal.Evaluation}, with the kernel
in \code{Internal.Measurement}. Only unitary actions differ.
\code{Measure ks} processes \code{reverse ks} and prepends each bit,
placing the new block in \code{ks} order before earlier outcomes.
\code{Initialize ks targets} measures those wires, applies $X$ wherever
the observed and target bits differ, and discards its new outcomes.
The RNG still advances. Initialization needs distinct, in-range wires and
matching target length; empty lists are no-ops. Reset returns a pure branch;
averaging over its hidden outcomes can produce a mixed state.
The stored squared norm need not be the probability of the recorded outcomes,
since resets hide samples.

\paragraph{Example.}
Using the Bell circuit of Section~\ref{sec:example} and $h=1/\sqrt2$:
\begin{lstlisting}
SV.evalProg [Unitary bell, Measure [0], Initialize [1] [False]]
            (St.ket [0,0]) [0.8,0.2]
\end{lstlisting}
returns a state with amplitudes $(0,0,h,0)$, outcomes \code{[True]}, and
an empty RNG. Measurement selected $h\ket{11}$; reset flipped its last bit.
The stored squared norm is $1/2$, but another measurement of wire zero
returns one with probability one, by \eqref{eq:measurement}.

\section{Resource growth and implementation evidence}
\subsection{Field growth: cost and ways to control it}
\label{sec:growth}
Three independent sizes grow: amplitudes or bonds, field degree $d$,
and rational coefficient bit lengths. Even product states can need large $d$.
For distinct odd prime angle denominators $p_1,\ldots,p_k$, the
sufficient-field rule can give
\[
N=8\prod_jp_j,\qquad d=\varphi(N)=4\prod_j(p_j-1).
\]
Rotations with denominators $101$ and $103$ give $N=83224$ and $d=40800$.
A ten-qubit statevector then has $2^{10}d=41\,779\,200$ rational-coordinate
slots before temporaries and integer payloads. Low entanglement alone
does not protect dense scalar coordinates. Polynomial addition costs
$O(d)$ and multiplication with reduction $O(d^2)$ rational operations;
integer growth adds further cost, hidden by the earlier field-operation counts.

Cancellation can leave a rational number in a huge ambient field.
Conversely, for distinct primes $p,q$,
$\zeta_p\zeta_q=\zeta_{pq}^{p+q}$ is primitive of order $pq$, since
$p+q$ is coprime to $pq$. Its degree really is $(p-1)(q-1)$.
Symbolic expressions can delay expansion but cannot remove this degree.
The following approaches offer partial remedies.

\paragraph{Reduce unnecessary field extensions.}
Use reduced angles and actual root orders:
$\zeta_{2b}^{a}$ has order $2b/\gcd(|a|,2b)$, and
$\zeta_{4b}^{a}$ has order $4b/\gcd(|a|,4b)$.
Validated rewrites can cancel inverses or combine same-axis rotations
before collecting roots. Preserve global phase: rotations have general
period $4$ in $q$, not $2$. Axis validation must precede cancellation.

\paragraph{Store fewer coordinates and reuse arithmetic.}
Sparse polynomials and special forms for rationals and roots of unity
avoid mostly empty arrays. Cache fields, reduction data, and embeddings.
Basis-state unitary evaluation stays in $D_N$, permitting integer
numerators with power-of-two denominators rather than general rational
arithmetic. Sparse products or reduction may nevertheless become dense,
and MPS elimination needs general division.

\paragraph{Delay promotion and expansion.}
Per-coefficient or per-core fields postpone promotion until values interact;
exact subfield descent can shrink results after cancellation. Both require
bookkeeping and checked conversions. Factored states or shared scalar
expressions can retain products: product-state cores may keep unrelated
angle fields separate until contraction, inner products, or equality.
Small bonds alone do not help if every core uses a large common field.
These optimizations are future work.

\subsection{Proof obligations and implementation evidence}
The \code{ExactTest} suite checks field invariants and backend agreement,
supplemented by independent expected values to detect shared errors:
$\sem{\code{R Z 1}}=-iZ$,
$\sem{\code{R (Id 0) 2}}=[-1]$,
$\sem{\code{Phase (1/2)}}=[i]$,
the worked examples, and QFTs through three qubits against the analytic DFT.
Further tests cover arithmetic, embeddings, composite operators, promotion,
rewrites, and failures. Randomized tests use exact equality and
shape-preserving shrinking. Run \code{cabal test all}.

\section{What exact MPS evaluation can and cannot mean}
\label{sec:mps}
\subsection{Representation and cut ranks}
A statevector has $2^n$ amplitudes even when its structure is simple.
For $n\ge1$, a matrix product state (MPS) stores those amplitudes as
products of small matrices:
\begin{equation}
v_{b_1\cdots b_n}
 =A^{[1]}_{b_1}A^{[2]}_{b_2}\cdots A^{[n]}_{b_n},
\qquad
A^{[j]}_b\in K^{\chi_{j-1}\times\chi_j},\quad \chi_0=\chi_n=1.
\label{eq:mps}
\end{equation}
The matrices $A^{[j]}_b$ are cores, $b\in\{0,1\}$ indexes the qubit,
and $\chi_j$ is a bond dimension. Site $j$ is QPP wire $j-1$.
Multiplication gives one scalar amplitude. Bonds at most $\chi$ require
$O(n\chi^2)$ storage. Product states have bonds one; entanglement can
require exponentially large bonds.

At width zero, store $z\in K$, with squared norm $\bar z z$.
\code{Phase q} already sends it to $e^{i\pi q}z$.
Tensoring this scalar with an MPS scales one core; two scalars multiply.
This base case includes zero, unlike an empty product of cores.

\paragraph{Unfolding a state across a cut.}
Group the first $k$ bits into a row index and the remaining bits into a
column index. This reshapes the amplitudes into a matrix
$M^{(k)}\in K^{2^k\times2^{n-k}}$.
An MPS factors this matrix through a space of dimension $\chi_k$,
so $\rank M^{(k)}\le\chi_k$. For a nonzero state, the minimum possible
bond dimension at each cut is this rank.
The Schmidt decomposition obtains the same rank from an SVD:
\[
M^{(k)}=U\Sigma V^\dagger,\qquad
\ket v=\sum_{\alpha=1}^{r}\sigma_\alpha
       \ket{\ell_\alpha}\otimes\ket{r_\alpha}.
\]
The left and right Schmidt vectors are orthonormal, and
$\sigma_\alpha^2$ are the nonzero eigenvalues of
$(M^{(k)})^\dagger M^{(k)}$ \cite{schollwock}.
The left vectors are columns of $U$; the right vectors are the
complex conjugates of columns of $V$.
Discarding small singular values gives controlled numerical approximation.

\subsection{The SVD obstruction}
Cyclotomic entries do not guarantee cyclotomic singular values.
For the integer matrix
\[
M=\begin{pmatrix}1&1&0\\0&1&1\\0&0&2\end{pmatrix},
\qquad \det(xI-M^TM)=x^3-8x^2+15x-4,
\]
the squared singular values are not cyclotomic. Consequently neither
are the singular values: squaring a cyclotomic number would stay
cyclotomic. The optional argument below proves this claim.

The squared entries of $M$ sum to $8$. Padding $M/\sqrt8$ to
$4\times4$ gives the amplitudes of a normalized four-qubit state whose
Schmidt values are noncyclotomic. It is QOp-preparable: a Hadamard and
two Hadamards controlled on its one branch give nonzero amplitudes
$(2,1,1,1,1)/\sqrt8$. A reversible computational-basis permutation
places them in the required entries. Such permutations can be built
from $X$ and multiply controlled $X$; they are more general than QPP's
wire-reordering \code{Permute}.

QPP's numerical \code{Exact} MPS setting only disables truncation;
its floating-point SVD does not provide algebraic exactness.

\paragraph{Optional proof of the obstruction (Galois theory).}
The displayed cubic has no rational root and discriminant
$916=4\cdot229$, a nonsquare. It is therefore irreducible with Galois
group $S_3$, so adjoining one root gives a non-Galois cubic field.
By the Galois criterion in Section~\ref{sec:normalization}, no root
of this cubic can belong to a cyclotomic field.
Dividing its positive square roots by $\sqrt8$, itself cyclotomic,
does not remove the obstruction for the normalized state.

\subsection{Rank factorization and exact local updates}
\label{sec:rank}
We need a factorization that removes linear dependence without finding
orthonormal bases. Let $T\in K^{m\times n}$ have rank $r>0$.
Choose $r$ independent columns with indices $J$ and form $B=T_{:,J}$.
Every column of $T$ is a unique linear combination of columns of $B$,
so $T=BC$ for some $C\in K^{r\times n}$.
Choose $r$ independent rows of $B$, indexed by $I$; then
\begin{equation}
T=BC,\qquad B=T_{:,J},\qquad C=(B_{I,:})^{-1}T_{I,:}.
\label{eq:rankfactor}
\end{equation}
The invertible $r\times r$ matrix $B_{I,:}$ makes the coefficients
computable using field arithmetic alone. For example,
\[
\begin{pmatrix}1&2&3\\2&4&6\end{pmatrix}
=\begin{pmatrix}1\\2\end{pmatrix}
 \begin{pmatrix}1&2&3\end{pmatrix};
\]
the column space has dimension one.

Implement the inverse action by elimination, not by forming an explicit
inverse. Select pivots by exact nonzero tests and handle the zero matrix
separately. The factors stay in $K$, although division may leave the
smaller dyadic ring $D_N$. Rank is unchanged on extending $K$ to $\C$:
it is the largest order of a nonzero minor, and those minors have the
same zero/nonzero status in either field.
Successively factor the unfoldings of a nonzero tensor to obtain an
exact MPS with minimal cut ranks; the rank-based existence argument
underlying tensor trains does not require Schmidt values \cite{oseledets}.

\paragraph{Applying a local gate.}
Contract two neighboring cores, apply the two-qubit gate matrix $U$,
and regard the result as a matrix:
\begin{equation}
T_{(\alpha,s),(t,\beta)}
=\sum_{a,b,\gamma}U_{st,ab}
 A^{[j]}_a(\alpha,\gamma)A^{[j+1]}_b(\gamma,\beta),
\qquad s,t,a,b\in\{0,1\}.
\label{eq:mpsupdate}
\end{equation}
The indices $\alpha,\gamma,\beta$ are bonds.
Factor $T=BC$ and reshape the factors into replacement cores.
With neighboring bonds at most $\chi$, this costs $O(\chi^3)$ field operations.
A one-qubit gate only mixes the two physical slices of one core.
Adjacent SWAPs can route a nonlocal gate to neighboring sites.

These cores need not be orthonormal. A change
$A^{[j]}_b\mapsto A^{[j]}_bG$ and
$A^{[j+1]}_b\mapsto G^{-1}A^{[j+1]}_b$ leaves amplitudes unchanged for
invertible $G$; this freedom is called a gauge.
Compute squared norms by contracting an environment through the cores:
\[
E_0=[1],\qquad
E_j=\sum_{b=0}^1 (A^{[j]}_b)^\dagger E_{j-1}A^{[j]}_b,\qquad
E_n=[\,v^\dagger v\,].
\]
One core alone generally does not contain the state's norm.

A local factorization need not give globally minimal bonds when surrounding
block bases are redundant; that requires rank-revealing sweeps.
No tolerance or truncation is involved. Bonds, degree, and coefficient
sizes remain independent costs; exactness guarantees no efficiency bound.

Hubig et al. study numerical MPO arithmetic with SVD, deparallelisation,
and delinearisation \cite{hubig}. The projects below study exact rank
decisions and coefficient growth in these representations.

\section{Student projects}
\titlespacing*{\paragraph}{0pt}{.25ex plus .25ex}{1em}
\subsection{HQP project: an exact MPS interpreter in four weeks}
\label{sec:hqp}
\paragraph{Starting point and scope.}
Supply scalars, oracles, tests, and a positioned-gate compiler before week 1;
the compiler is still a supervisor task. Prerequisites: Haskell, linear algebra,
and circuits.
Implement checked cores, basis states, contraction, local updates, and rank
factorization. Support $X,Y,Z,H,\mathrm{SX}$, rotations about $X,Y,Z$, global phase,
CNOT, controlled rational phase, and SWAP; all two-qubit gates are adjacent.
The supplied compiler handles composition, tensor padding, and adjoints in
this fragment. Reject other control forms, \code{DirectSum}, composite axes, and
nonadjacent gates.

\paragraph{Weekly milestones.}
Week 1: deterministic rank factorization; verify $T=BC$ and ranks on
rectangular, deficient, and zero matrices.
Week 2: cores, contraction, and updates; compare Bell/GHZ and inverse circuits
with the oracle after every gate.
Week 3: gate interface, adjacent SWAPs, environment squared norms, random
circuits, and errors. Week 4: experiments, report, and demonstration.

\paragraph{Acceptance and deliverables.}
Deliver code, tests, benchmarks, and a 6--10-page report. Up to eight qubits,
check amplitudes, squared norms, and inverse round trips, including cancellation,
non-dyadic angles, controlled phases, and SWAP ordering. Vary depth and
denominators on product, GHZ, and shallow entangling circuits; report time,
memory, bonds, degree, coefficient bits, and failures. No speedup is required.
Local rank reduction is mandatory; global sweeps and routing are stretch work.
General \code{QOp}, measurement, SVD, and approximation are excluded.

\subsection{MSc thesis: exact tensor-network semantics and coefficient growth}
\label{sec:msc}
\paragraph{Questions and required baseline.}
Measure MPS/MPO time and memory savings, effects of pivot/compression policies,
and usefulness for compiler validation. MPO cores have input/output bits and two bonds.
Build a complete compositional \code{QOp} interpreter in one field, without
dense fallback. Prove gauge, wire-order, and semantic invariants, including
bidirectional sweeps that remove redundant block bases. Implement control
and direct sum by MPO arithmetic or direct branching, using Section~\ref{sec:matrix};
rotations use \eqref{eq:svrotation}. Addition concatenates or block-diagonalizes
bonds before rank compression. Handle zero branches.

\paragraph{Experimental scope.}
Choose one comparison: first-nonzero versus coefficient-size pivots, or eager
versus bond-triggered compression. A second comparison is stretch work.
Local fields require separate promotion, contraction, and equality invariants;
fraction-free and modular elimination are optional.

\paragraph{Six-month milestones.}
Month 1: literature, prototype, design. Months 2--3: baseline, sweep proofs,
exact-agreement tests. Month 4: comparison and instrumentation. Month 5:
experiments and rewrites. Month 6: thesis and reproducible release.

\paragraph{Evaluation.}
Compare exact Matrix/Statevector, uncompressed MPS, and floating MPS.
Vary bonds and degrees independently on product states, QFTs, Clifford+$T$,
mixed-denominator circuits, and rewrite regressions. Publish seeds, limits,
builds, results, and timeouts. Success requires correctness and explained costs;
measurement and parallelism are excluded.

\subsection{Four-week project: efficient cyclotomic arithmetic}
\label{sec:scalarproject}
Preserve the scalar API. Week 1: tests and instrumentation; week 2: sparse
polynomials; week 3: sparse/dense switching and cached fields/embeddings;
week 4: evaluation and report. Exclude minimal-field discovery.
Deliver code, tests, benchmarks, and a report. Compare arithmetic, conjugation,
inversion, promotion, equality, and ordering against the prototype and, where
supported, Walck \cite{walck}. Test random values, cancellation, circuits, and
elimination; report time, memory, degree, sparsity, coefficient bits, and failures.

\begin{thebibliography}{9}
\enlargethispage{12pt}
\footnotesize
\setlength{\itemsep}{0pt}
\setlength{\parsep}{0pt}
\setlength{\parskip}{0pt}
\bibitem{milne}
J.~S. Milne. \href{https://www.jmilne.org/math/CourseNotes/FT.pdf}
{\emph{Fields and Galois Theory}}, version 5.10, 2022.

\bibitem{walck}
S.~N. Walck. \href{https://hackage.haskell.org/package/cyclotomic-1.1.2}
{\emph{cyclotomic: A subfield of the complex numbers for exact calculation}},
Haskell package 1.1.2, 2023.

\bibitem{base}
GHC libraries. \href{https://hackage.haskell.org/package/base-4.21.0.0/docs/Data-Complex.html}
{\emph{Data.Complex}}, \code{base-4.21.0.0} documentation.

\bibitem{giles}
B. Giles and P. Selinger. \href{https://doi.org/10.1103/PhysRevA.87.032332}
{Exact synthesis of multiqubit Clifford+$T$ circuits}.
\emph{Physical Review A} 87, 032332, 2013.

\bibitem{schollwock}
U. Schollw\"ock. \href{https://doi.org/10.1016/j.aop.2010.09.012}
{The density-matrix renormalization group in the age of matrix product states}.
\emph{Annals of Physics} 326, 96--192, 2011.

\bibitem{oseledets}
I.~V. Oseledets. \href{https://doi.org/10.1137/090752286}{Tensor-train decomposition}.
\emph{SIAM Journal on Scientific Computing} 33(5), 2295--2317, 2011.

\bibitem{hubig}
C. Hubig, I.~P. McCulloch, and U. Schollw\"ock.
\href{https://doi.org/10.1103/PhysRevB.95.035129}
{Generic construction of efficient matrix product operators}.
\emph{Physical Review B} 95, 035129, 2017.
\end{thebibliography}
\end{document}
